\section{Introducción}

El crecimiento de los asistentes de voz y los sistemas domóticos ha
popularizado el uso de modelos de lenguaje para interpretar comandos en
lenguaje natural (``prendé la luz del living'', ``bajá el aire a 20
grados''). La mayoría de las soluciones comerciales dependen de modelos de
gran escala alojados en la nube, lo que introduce tres problemas recurrentes
en el contexto de un hogar: latencia de red, costo recurrente por uso de
API, y dependencia de conectividad — un comando tan simple como apagar una
luz no debería fallar por una caída de internet.

Una alternativa creciente es el uso de Small Language Models (SLMs) —
modelos de lenguaje de entre unos cientos de millones y pocos miles de
millones de parámetros — ejecutados localmente en el propio hub domótico
(una Raspberry Pi, un mini-PC o un microcontrolador de gama alta). Trabajos
recientes muestran que estos modelos pueden alcanzar, en tareas acotadas, un
desempeño competitivo frente a modelos mucho más grandes \cite{gupta2025slms}.
Sin embargo, la mayor parte de la literatura sobre SLMs y asistentes
domóticos evalúa el inglés, y no está claro cómo se comportan estos modelos
frente a estructuras propias del español rioplatense (voseo, diminutivos,
expresiones relativas como ``bajale un toque'') ni cuál es el costo real en
latencia de ejecutarlos en hardware modesto.

La clasificación manual de respuestas incorrectas para identificar patrones
de error no escala a un roster de doce modelos: multiplica el volumen de
respuestas a revisar a mano, y su resultado tampoco es reproducible por un
tercero, porque depende del criterio del anotador en el momento de
clasificar cada caso. Este trabajo evalúa doce SLMs mediante un pipeline de
verificación completamente automático, sobre un roster que incorpora
familias de modelos publicadas recientemente (LFM2.5, Granite 4.0, Qwen3.5
y OLMo-2) junto con dos tiers de tamaño (menos de 1000 millones y entre
1000 y 2000 millones de parámetros).

Las contribuciones de este trabajo son: (1) la evaluación de un roster de
doce modelos abiertos sub-2B, de acceso público y sin restricción de
licencia ni de token de acceso, organizados en dos tiers de tamaño de seis
modelos cada uno y evaluados bajo un protocolo común; (2) un pipeline de
verificación automática en dos etapas, coincidencia textual determinista
más un juez de lenguaje con una taxonomía cerrada, que elimina el
etiquetado manual de respuestas incorrectas y de categorías lingüísticas;
(3) el reporte conjunto de exactitud estricta y exactitud laxa, cuya
brecha cuantifica el costo de exigir coincidencia textual exacta; (4) un
protocolo de ejecución reproducible, con aislamiento por contenedor y
versiones de librerías documentadas por grupo; y (5) la inclusión
deliberada de un modelo representativo de la generación anterior de SLMs
(\texttt{Qwen2.5-1.5B-Instruct}), evaluado bajo el mismo protocolo que el
resto del roster, que permite cuantificar cuánto cambió lo factible para
esta tarea en el rango sub-2B respecto de esa generación.
